\documentclass[../DoAn.tex]{subfiles}
\begin{document}
\section{Kết luận}
Trải qua quá trình nghiên cứu, phân tích và phát triển, đồ án đã hoàn thành mục tiêu xây dựng "Hệ thống tổ chức thi trắc nghiệm và giám sát thi cử trực tuyến áp dụng tại Đại học Bách Khoa". So với các nền tảng truyền thống thường phân mảnh giữa học và thi, JQK Study đã tích hợp thành công hai luồng nghiệp vụ này vào chung một hệ thống, mang lại trải nghiệm liền mạch cho người dùng.

Những đóng góp nổi bật của hệ thống bao gồm:
\begin{itemize}
    \item \textbf{Tự động hóa ra đề:} Giúp giảng viên tiết kiệm thời gian đáng kể thông qua thuật toán sinh đề thi ngẫu nhiên dựa trên cấu hình chủ đề và độ khó.
    \item \textbf{Giám sát thi cử hiệu quả:} Xây dựng được cơ chế chống gian lận (Anti-Cheat) theo thời gian thực bằng việc bắt các sự kiện chuyển tab, copy/paste, đảm bảo tính công bằng của kỳ thi.
    \item \textbf{Quản lý khép kín:} Mô hình tổ chức Lớp học - Khóa học - Đề thi giúp quản lý tiến độ học tập và điểm số một cách thống nhất, tập trung.
\end{itemize}

Bên cạnh những kết quả đạt được, hệ thống vẫn còn một số hạn chế nhất định như chưa hỗ trợ đa dạng hơn các loại hình câu hỏi tự luận hay câu hỏi có tính tương tác cao, cơ chế giám sát chưa áp dụng công nghệ nhận diện khuôn mặt qua AI để tăng cường bảo mật.

\section{Hướng phát triển}
Trong tương lai, để hoàn thiện và nâng cao năng lực của hệ thống JQK Study, các hướng phát triển tiếp theo bao gồm:
\begin{itemize}
    \item \textbf{Tích hợp AI trong giám sát (Proctoring):} Ứng dụng trí tuệ nhân tạo để phân tích hình ảnh từ webcam, nhận diện khuôn mặt và phát hiện các hành vi bất thường của người thi.
    \item \textbf{Đa dạng hóa câu hỏi:} Mở rộng hệ thống để hỗ trợ câu hỏi tự luận, bài tập lập trình (chấm code tự động) nhằm phục vụ các ngành học đặc thù.
    \item \textbf{Phân tích dữ liệu học tập (Learning Analytics):} Áp dụng các kỹ thuật khai phá dữ liệu để phân tích kết quả bài thi, từ đó đưa ra các gợi ý cá nhân hóa nhằm cải thiện lộ trình học tập cho từng học viên.
\end{itemize}

\end{document}